\documentclass[10pt,a4paper]{amsart}
\usepackage[margin=19mm]{geometry}
\usepackage{amsmath,amssymb,mathtools}
\usepackage[scr=rsfs,cal=euler]{mathalfa}
\usepackage{xfrac}
\usepackage[unicode,hidelinks,bookmarks=false]{hyperref}
\newcommand{\ie}{i.e.\ }
\newcommand{\cf}{cf.\ }
\newcommand{\resp}{respectively\ }
\newcommand{\Subsection}{\subsection}
\providecommand{\nobreakdash}{\nobreak\hbox{-}}
\setlength{\emergencystretch}{2em}
\multlinegap0pt
\usepackage{amssymb}
\usepackage{mathtools}
\usepackage[shortlabels]{enumitem}
\usepackage{float}
\newtheorem{prop}{Proposition}
\title{How to Study Reflected Brownian Motion in a Quadrant via Kernel Functional Equations?\\ \medskip \large A short survey}
\author{Sandro Franceschi}
\date{Source: arXiv:2606.06187v1 (CC BY 4.0)}
\begin{document}
\maketitle

\noindent\emph{Source and licence.} How to Study Reflected Brownian Motion in a Quadrant via Kernel Functional Equations?\\ \medskip \large A short survey. Version arXiv:2606.06187v1 (CC BY 4.0), distributed by arXiv under the Creative Commons Attribution 4.0 International licence, http://creativecommons.org/licenses/by/4.0/, as recorded in the arXiv licence metadata for this exact version and shown on the abstract page https://arxiv.org/abs/2606.06187v1. Full licence notice: \texttt{licenses/arxiv-2606.06187-CC-BY-4.0.txt}.\par\smallskip

\noindent\textit{Editorial scope.} Counted unit: the article's prop eq:Ac (source0.tex lines 2921--2994). The article's complete original proof of it is delivered with it (source0.tex lines 2995--3055, one \texttt{proof} environment of 61 lines). No mathematics is written, repaired or re-derived: the statement and the proof are the article's own lines, in the article's own notation, and no retained line is substituted or re-flowed.  No further source-local context is needed: every cross-reference of every retained line resolves inside the two delivered spans.  Nothing was omitted from any retained span, so the package carries no omission record.  No retained line contains a citation, so the wrapper carries no bibliography and no bibliography entry.  No retained line contains a figure environment, an image-inclusion command or drawing code, so no image is delivered and the package declares no asset dependency.  Editorial formatting only, in addition: an AMS \texttt{amsart} preamble that restates the article's own declarations and macros used by the retained lines, namely the environments \texttt{prop} on separate counters, together with the standard packages those lines need.  The retained environments print in this document as Proposition 1 in order of appearance, whereas the article prints the same items with the numbering of its own full document.  The complete native source of the article as distributed in the arXiv e-print for this exact version is bundled unchanged as \texttt{sources/202167/source0.tex}.
\begin{prop}[Canonical normalization of the reflected Brownian motion]
Let $(Z_t)_{t\ge 0}$ be a reflected Brownian motion in $\mathbb{R}_+^2$ satisfying
\[
Z_t = Z_0 + \mu t + B_t + R L_t,
\]
where the covariance matrix, the reflecting matrix and the drift are
\[
\Sigma =
\begin{pmatrix}
\sigma_{11} & \sigma_{12}\\
\sigma_{12} & \sigma_{22}
\end{pmatrix},
\qquad
R =
\begin{pmatrix}
r_{11} & r_{12}\\
r_{21} & r_{22}
\end{pmatrix},
\qquad
\mu= \begin{pmatrix}
\mu_1 \\
\mu_2 
\end{pmatrix}.
\]
Define
\begin{equation}
A =
\begin{pmatrix}
\sigma_{11}^{-1/2} & 0\\
0 & \sigma_{22}^{-1/2}
\end{pmatrix},
\qquad
c = |A\mu|= \sqrt{
\frac{\mu_1^2}{\sigma_{11}} + \frac{\mu_2^2}{\sigma_{22}}
},
\label{eq:Ac}
\end{equation}
and introduce the process
\[
\bar Z_t = c\, A\, Z_{t/c^2}.
\]
Then $(\bar Z_t)_{t\ge 0}$ is a reflected Brownian motion in $\mathbb{R}_+^2$ satisfying
\[
\bar Z_t = \bar Z_0 + \bar\mu t + \bar B_t + \bar R\, \bar L_t,
\]
where the covariance matrix, the reflecting matrix and the drift are
\[
\bar\Sigma =
\begin{pmatrix}
1 & \rho\\
\rho & 1
\end{pmatrix}
\text{ with }
\rho = \frac{\sigma_{12}}{\sqrt{\sigma_{11}\sigma_{22}}},
\quad
\bar R =
\begin{pmatrix}
1 &
\sqrt{\frac{\sigma_{22}}{\sigma_{11}}} \frac{ r_{12}}{r_{22}}\\[1em]
\sqrt{\frac{\sigma_{11}}{\sigma_{22}}} \frac{ r_{21}}{r_{11}} &
1
\end{pmatrix},
\quad
\bar\mu = \frac{A\mu}{|A\mu|}
\text{ with }
|\bar\mu|=1,
\]
and
\[
\bar L_t^i = \frac{cr_{ii} }{\sqrt{\sigma_{ii}}}\; L_{t/c^2}^i,
\qquad i=1,2,
\]
is the local time of the process on the axes.
\end{prop}

\begin{proof}
Set
\[
A=\begin{pmatrix}
\sigma_{11}^{-1/2} & 0\\
0 & \sigma_{22}^{-1/2}
\end{pmatrix},
\qquad
c=|A\mu|,
\qquad
D=\begin{pmatrix}
r_{11} & 0\\
0 & r_{22}
\end{pmatrix}.
\]
By definition,
\[
\bar Z_t
=cA Z_{t/c^2}=
cAZ_0+\frac{A\mu}{c}\,t+cA B_{t/c^2}+cA R L_{t/c^2}.
\]
By Brownian scaling, \(\bar B_t := cA B_{t/c^2}\) is a Brownian motion with covariance matrix
\[
A\Sigma A^\top
=
\begin{pmatrix}
1 & \rho\\
\rho & 1
\end{pmatrix},
\quad\text{where}\quad
\rho=\frac{\sigma_{12}}{\sqrt{\sigma_{11}\sigma_{22}}}.
\]
The drift becomes
\[
\bar\mu=\frac{A\mu}{|A\mu|}, \qquad |\bar\mu|=1.
\]
To conclude, it suffices to observe that
\[
cA R L_{t/c^2}
=
\underbrace{A R A^{-1}D^{-1}}_{=\bar R}
\underbrace{c D A L_{t/c^2}}_{=\bar L_t}.
\]
The rescaled local time vector is
\[
\bar L_t
=
cD A\, L_{t/c^2},
\]
that is,
\[
\bar L_t^i=\frac{c\, r_{ii}}{\sqrt{\sigma_{ii}}}\,L_{t/c^2}^i,
\qquad i=1,2.
\]
The new reflection matrix is
\[
\bar R=AR A^{-1}D^{-1}.
\]
A direct computation yields the expression of $\bar R$ given in the proposition.
This proves that \((\bar Z_t)_{t\ge0}\) is a reflected Brownian motion in \(\mathbb{R}_+^2\) with parameters \(\bar\Sigma\), \(\bar R\), \(\bar\mu\), and local time \(\bar L_t\).
\end{proof}

\end{document}
